\begin{tikzpicture}[font=\small,>=stealth]
\draw[rounded corners=3pt,gray!45] (0,0) rectangle (14.2,11.3);
\node at (7.1,10.75) {先认条件：通径、普通竖弦、轴线，只有第一幅能直接用定理};
\node at (2.55,9.75) {过焦点的通径};
\begin{scope}[shift={(1.9,5.5)},scale=0.65]
\draw[->,gray] (-1.2,0)--(4.3,0) node[above] {$x$};
\draw[teal!75!black,thick,samples=120,domain=-3.85:3.85,variable=\t] plot ({\t*\t/4},{\t});
\draw[blue!75!black,very thick] (1,0)--(1.0000000000000002,2.0000000000000004);
\draw[orange!85!black,very thick] (1,0)--(0.9999999999999999,-2.0);
\fill (1.0000000000000002,2.0000000000000004) circle(2pt) node[above=8pt] {$A$};
\fill (0.9999999999999999,-2.0) circle(2pt) node[below left=6pt] {$B$};
\fill (1,0) circle(2pt) node[below right=5pt] {$F$};
\end{scope}
\node at (2.55,1.95) {$1/2+1/2=1$};
\node at (6.95,9.75) {不经过焦点的普通弦};
\begin{scope}[shift={(6.3,5.5)},scale=0.65]
\draw[->,gray] (-1.2,0)--(4.3,0) node[above] {$x$};
\draw[teal!75!black,thick,samples=120,domain=-3.85:3.85,variable=\t] plot ({\t*\t/4},{\t});
\draw[gray,dashed] (2,-2.8284271247461903)--(1,0)--(2,2.8284271247461903);
\draw[blue!75!black,very thick] (2,-2.8284271247461903)--(2,2.8284271247461903);
\fill (2,-2.8284271247461903) circle(2pt);\fill (2,2.8284271247461903) circle(2pt);
\fill (1,0) circle(2pt) node[above=8pt] {$F$};
\end{scope}
\node at (6.95,1.95) {$1/3+1/3=2/3$};
\node at (11.35,9.75) {轴线：只有一个交点};
\begin{scope}[shift={(10.7,5.5)},scale=0.65]
\draw[->,gray] (-1.2,0)--(4.3,0) node[above] {$x$};
\draw[teal!75!black,thick,samples=120,domain=-3.85:3.85,variable=\t] plot ({\t*\t/4},{\t});
\draw[red!75!black,very thick] (-.6,0)--(4.1,0);\fill (0,0) circle(2pt) node[below left=5pt] {$O$};
\fill (1,0) circle(2pt) node[below right=5pt] {$F$};
\end{scope}
\node at (11.35,1.95) {不构成有限两点弦};

\node at (7.1,.7) {三图同为 $y^2=4x$、$p=2$；普通弦的两端到 $F$ 都是 $3$，不满足倒数和 $1$。};
\end{tikzpicture}
